% =====================================================================
% 3. EXPLORATORY DATA ANALYSIS (EDA)
% =====================================================================
\section{Exploratory Data Analysis (EDA)}\label{sec:prdl-eda}

\subsection{Introduction}

The EDA has a concrete purpose in a supervised problem: to find which variables separate the
classes or relate to the target, to detect problems (leakage, redundancy, rare categories) before
modelling, and to set realistic expectations. Its structure follows the course EDA notebook. Because
almost every predictor here is binary, the techniques were adapted, and an unsupervised analysis of
the mutation profiles was added. All EDA statistics describe the full dataset; none of them is used
to select predictors, because selecting variables with the labels outside the validation loop biases
the error estimate \cite{ambroise2002}.

\subsection{Techniques}

\textbf{Tests and effect sizes.} Age is not normally distributed within either class (Shapiro--Wilk
test), so the two classes are compared with the Mann--Whitney $U$ test \cite{mann1947} rather than
a $t$-test. Its statistic has a direct link with classification: $U/(n_1 n_0)$ is the probability
that a random GBM patient is older than a random LGG patient, which is the area under the ROC curve
(AUC) of age used alone \cite{hanley1982}. Binary predictors are compared with Fisher's exact test
\cite{fisher1922}, which, unlike the $\chi^2$ test, remains valid when some cells of the $2\times2$
table are small, as happens for rarely mutated genes. Since a $p$-value does not measure the size of
an effect, every test is reported with an effect size (Cohen's $d$, odds ratio, single-feature AUC
with a bootstrap interval \cite{efron1993}). With \NTests{} simultaneous tests, $p$-values are
corrected with Bonferroni (control of any false positive) and Benjamini--Hochberg (control of the
expected proportion of false discoveries, less conservative) \cite{benjamini1995}.

\textbf{Redundancy.} For binary variables the Pearson correlation is the phi coefficient. The
variance inflation factor, $\mathrm{VIF}_j = 1/(1-R_j^2)$, measures how well a predictor is
explained by all the others; values above 5--10 signal problematic collinearity \cite{james2021}.

\textbf{Unsupervised techniques.} PCA finds orthogonal directions of maximum variance after
centring the data; standardising 0/1 variables gives a rare mutation the same weight as a frequent
one, so both versions are compared. Euclidean geometry treats two patients as similar when they
share the \emph{absence} of mutations, which dominates sparse data. The Jaccard distance,
$d_J(A,B) = 1 - |A\cap B|/|A\cup B|$, ignores shared absences; it is used for metric MDS, whose
fit is measured with Kruskal's stress \cite{kruskal1964}, and for t-SNE, a non-linear map that
preserves local neighbourhoods \cite{vandermaaten2008}. $k$-means minimises the within-cluster
variation; on 0/1 data the squared Euclidean distance equals the number of genes in which two
patients differ. Hierarchical clustering needs a dissimilarity and a linkage; Ward's linkage merges
the groups that least increase the within-cluster variance \cite{ward1963}. Partitions are compared
with the silhouette \cite{rousseeuw1987} and the adjusted Rand index (ARI), which is 0 for chance
agreement and 1 for identical partitions \cite{hubert1985}.

\subsection{Results \& Conclusions}

\textbf{The classes differ strongly in a few variables.} GBM patients are older (median
\AgeMedianGBM{} vs \AgeMedianLGG{} years; Cohen's $d = \AgeCohenD$; AUC \AUCAge{} \AUCAgeCI), and the
largest single difference is IDH1, mutated in \PrevIDHOneLGG\% of LGG and \PrevIDHOneGBM\% of GBM
patients (AUC \AUCIDHOne{} \AUCIDHOneCI; Figure~\ref{fig:prdl-prevalence}). PTEN and EGFR mutations
are more frequent in GBM, and CIC and NOTCH1 almost exclusive to LGG (NOTCH1:
\PrevNOTCHLGG\% vs \PrevNOTCHGBM\%), which anticipates the separation problems of logistic
regression (Section~\ref{sec:prdl-results}). After correction, \NSigBonf{} of \NTests{} differences
remain significant with Bonferroni and \NSigBH{} with Benjamini--Hochberg; most of the remaining
genes have AUCs close to 0.5, so statistical significance does not imply usefulness.

\begin{figure}[H]
\centering
\includegraphics[width=0.62\textwidth]{eda_mutation_prevalence.png}
\caption{Percentage of patients with each mutation in LGG and GBM (sorted by the difference).}
\label{fig:prdl-prevalence}
\end{figure}

\textbf{The redundancy is moderate and biologically structured.} The strongest correlations are
age--IDH1 ($r = \CorrAgeIDH$), TP53--ATRX ($r = \CorrTPATRX$) and CIC--FUBP1 ($r = \CorrCICFUBP$),
and the maximum VIF is \MaxVIFEDA, below the usual thresholds. Correlated mutations do not make linear
models unstable here, but they violate the independence assumption of naive Bayes.

\textbf{Leakage would make the problem look solved.} A quick cross-validated logistic regression
reaches an AUC of \BaseAUCAll{} with the 23 legitimate predictors and \BaseAUCAgeIDH{} with only age
and IDH1, but \BaseAUCLeak{} when the histological diagnosis is added. The last number measures
the leakage, not the information in the data.

\textbf{The unsupervised structure is the IDH axis.} PCA does not compress the profiles: \NCompCentred{}
centred (or \NCompScaled{} scaled) components are needed to explain 80\% of the variance, and PC1,
which explains \PCOneCentred\% of the centred variance, is the contrast between IDH1/TP53/ATRX and
PTEN/EGFR. With the Jaccard distance, \PctJaccardOne\% of the patient pairs share no mutation at all
($d_J = 1$), and the MDS map distorts the distances considerably (stress \MDSStress). Hierarchical
clustering with Jaccard distance fails: the \NNoMutation{} patients without any mutation are at
distance 1 from everybody, so complete linkage ties at 1 and merges all patients into one cluster,
while average and single linkage only isolate those \NNoMutation{} patients. This is a property of the
distance, not a bug, and it illustrates why the dissimilarity must be chosen for the data.
$k$-means and Ward's method on the 0/1 profiles agree well. The number of clusters was fixed before
looking at the result by a rule (highest agreement between the two algorithms), which gave
$k = \KMain$ (ARI \ARIKMWardMain). The two clusters are essentially IDH-wildtype
(\CluMainAN{} patients, \CluMainAGBM\% GBM, median age \CluMainAAge) and IDH-mutant (\CluMainBN{} patients,
\CluMainBGBM\% GBM, median age \CluMainBAge), with ARI \ARIGradeMain{} against the grade. As an
exploratory, secondary solution, $k = \KSecond$ splits the IDH-mutant group into a TP53/ATRX profile
and a CIC/FUBP1 profile with \CluSecondCGBM\% GBM (Figure~\ref{fig:prdl-clusters}), a structure that
mirrors the TCGA description of lower-grade gliomas: IDH-mutant tumours without 1p/19q codeletion
nearly all had TP53 mutations (94\%) and ATRX inactivation (86\%), while codeleted ones carried CIC,
FUBP1 and NOTCH1 mutations \cite{tcga2015lgg}. The clusters thus recover, from point mutations only,
the astrocytoma/oligodendroglioma distinction that the classification makes with the codeletion. Because the
cluster structure was not imposed by the grade, its agreement with the grade supports H3, and its
lack of perfect agreement is consistent with the outdated labels discussed in Section 2.1.

\begin{figure}[H]
\centering
\includegraphics[width=0.85\textwidth]{unsup_cluster_profiles.png}
\caption{Mutation prevalence in each $k$-means cluster for the main ($k=\KMain$) and the secondary
($k=\KSecond$) solution. The grade was not used to build the clusters.}
\label{fig:prdl-clusters}
\end{figure}

\textbf{Conclusions for modelling.} (i) Accuracy alone is misleading because the majority class
already gives \MajorityAcc; AUC, sensitivity and specificity are reported. (ii) Age and IDH1 carry most
of the signal, so large gains from complex models are not expected. (iii) Rare categories (one patient
of one race, mutations present in a single class) require penalised models and encoders that tolerate
unseen categories. (iv) The diagnosis and the project must never be predictors.
